\section{Choix du model pour la musique et la guitare}

\subsection{Modèle pour les accords}

Nous devons dans un premier temps modéliser l'écriture des accords.

\begin{defi}[Note de musique]
  Une note de musique permet de décrire la hauteur et la durée d'un son. Cependant en
  fonction du tempérament de la gamme utilisée, les fréquences changent.
  Par soucis de généralités, on va utiliser la convention MIDI \cite{bMIDI} 
  pour numéroter les notes.\\
  Soit :
  \begin{itemize}
    \item Le do au centre du piano (do3) est numéroté par 60.
    \item On numérote relativement à do3 les autres notes en ajoutant/retirant 1 par
      demi ton.
  \end{itemize}
\end{defi}

\begin{defi}[Accord]
  Un accord est un ensemble de notes.
  Un accord est donné par sa fondamentale (note de base pour la construction),
  et de notes souvent la tierce, la quinte, parfois les septièmes, neufièmes et
  treizièmes. Ainsi que la manière de positionner ces notes.\\
  Il existe pour cela différentes manières de noter les accords.
\end{defi}

\begin{defi}[Notation Américaine]
  On donne la fondamentale avec une lettre $[A\rightarrow G]$ allant de La
  à Sol que l'on peut coupler avec une altération \musFlat{}, \musSharp{}.\\
  On ajoute une indication sur la tierce \texttt{m} pour mineur et \texttt{M}
  pour majeur, on peut aussi indiquer \texttt{sus2} (resp. \texttt{sus4}) dans
  le cas où l'on joue la seconde (resp. la quarte) à la place de la tierce.\\
  On ajoute une indication pour savoir si la septième est présente avec
  \texttt{7}, de même pour la neuvième et la treizième.\\
  Il y a des cas exeptionels :
  \begin{itemize}
    \item Accord diminué \texttt{dim}
    \item Accord avec la sixte \texttt{6}
    \item Accord avec la neuvième sans la septième \texttt{add9}
    \item Accord de dominante \texttt{7}
  \end{itemize}
  Cette notation est presque identifiable à une grammaire :
  $$Acc := \left[A\rightarrow
  G\right]+\left[\musFlat{}|\musSharp{}|\varepsilon\right]+\left[m|M|sus2|sus4|\varepsilon\right]+\left[5-|5+|\varepsilon\right]+\left[6|\varepsilon\right]+\left[7m|7M|\varepsilon\right]+\left[9|\varepsilon\right]+\left[11|\varepsilon\right]+\left[13|\varepsilon\right]$$
\end{defi}



\begin{defi}[Choix de la notation]
  Nous décidons d'opter pour une notation de l'accord comme suit :
  \begin{itemize}
    \item Un nombre pour la note de construction de l'accord
    \item Une suite de nombre corrsepondant au décalage en demi tons à la
      première note.
  \end{itemize}
  Alors l'accord de do majeur pourra s'écrire ainsi :
  $$(60, 0, 4, 7)$$
\end{defi}

\begin{rema}
  Les intérêts d'une telle notation :
  \begin{itemize}
    \item Très simple d'utilisation
    \item On peut convertir les différents systèmes de chiffrages à celui-ci.
    \item Les autres chiffrages ne correspondent pas à des grammaires précises.
  \end{itemize}
  Il y a cependant des inconvenients :
  \begin{itemize}
    \item On ne peut pas différencier les positions dans l'accord (ce que l'on
      peut faire avec la notation américaine)
    \item Ainsi si l'on shouaite éliminer des notes (comme la fondamentale et/ou
      la quinte, intéréssant si la guitare n'est pas le seul instrument) on ne
      peut pas le faire aisément.\\
      Exemple : \texttt{C13} :
      $$(60, 0, 4, 7, 10, 14, 17) \textbf{ ou } (60, 0, 2, 4, 5, 7, 10)$$
      quelle note est la quinte ?
  \end{itemize}

  Nous avons donc ici un problème qui pourra resortir au moment où l'on se
  demandera comment faire pour éliminer judicieusement les notes dans un
  accord.
\end{rema}



\subsection{Modélisation de la guitare}

On définit une guitare par :
\begin{itemize}
  \item \texttt{nc} : nombre de cordes (int)
  \item \texttt{nf} : nombre de frêtes (int)
  \item \texttt{sc} : scordature de la guitare (vector[ ])
\end{itemize}

On définit un guitariste par :
\begin{itemize}
  \item \texttt{nfmi} : la taille de la main en le nombre de frêtes accessibles
    depuis le haut du manche (int)
	\item \texttt{nbdo} : le nombre de doigts utilisés pour jouer (int)
	\item \texttt{isba} : s'il est possible de réaliser un barré (bool)
  \item \texttt{ba} : le nombre de frêtes accessible avec un barré (-1 si barré
    impossible, 0 si barré et c'est tout, sinon on peut appuyer à d'autres
    endroits)  (int)
  \item \texttt{po} : le nombre de cordes bloquables avec le pouce (0 pouce
    inutilisé, 1 ou 2 nb cordes) équivaut à un petit barré (int)
\end{itemize}

\section{Le model de contraintes}

On définit les variables indéterminées 
$$D_1,D_2,D_3,D_4,D_5\in\left\{\left\{1, \dots, nf\right\}\times
\left\{1, \dots,6\right\}\right\}\cup\left\{\varepsilon\right\}$$
qui décrivent la position des doigts sur le manche.

Cependant il y a de nombreuses solutions à éliminer comme avoir plusieurs
doigts sur une même corde.
Plus généralement, l'espace des solutions potentielles est très vaste. Et l'on
a trouver une modélisation qui est plus simple à utiliser avec un plus petit
espace de solutions.

\subsection{Un second model de contraintes}

On définit les variables indéterminées $\left\{F_i\right\}_{1\leq i\leq nc}$ de
domaine $\left\{-1, 0, 1, \dots, nf\right\}$, on ne représente ici que
l'endroit où chaque corde est appuyée et non les doigts du guitarise.\\
-1 : corde non jouée, 0 : corde à vide, reste : numéro de la frête utilisée.

Premièrement, il faut s'assurer que l'on joue exactement les notes de l'accord, soit que :
$$\left\{\text{Notes jouées}\right\} \subset \left\{\text{Notes de
l'accord}\right\}$$

Il faut définir une écart entre les doigts de la sorte :
$$ecart = f(\max_n F_n, \min_n F_n)$$

Première contrainte :\\
Il faut que les notes jouées soient incluses dans l'accord, ie :
\begin{equation}
\forall i\in \llbracket1,nc\rrbracket\exists j\geq 1, (sc[i]+F_i)-(accord[0]+accord[j])
= 0(mod 12)
\end{equation}
Il faut que les notes de l'accord soient incluses dans les notes jouées, ie :
\begin{equation}
\forall i\in \llbracket1,na\rrbracket\exists j\geq 1, (sc[j]+F_j)-(accord[0]+accord[i])
= 0(mod 12)
\end{equation}
Il faut que les notes de l'accord soient incluses dans les notes jouées, ie :
cest contraintes nous assurent que l'on joue toutes les notes de l'accord (il
faudra par la suite modifier la deuxième contrainte.

Deuxième contrainte :\\
Il ne faut pas que l'écart soit trop important, et ce fixé par rapport
à \texttt{tm}.

